\documentclass[12pt]{beamer}

\usepackage[utf8]{inputenc}
\usepackage[T5]{fontenc}
\usepackage[vietnamese]{babel}
\usepackage{lmodern} % Đã có font hỗ trợ tốt tiếng Việt
\usepackage{csquotes} % Thêm gói này để biblatex không báo lỗi
% \setbeamertemplate{bibliography item}{\insertbiblabel}
\usepackage{amsmath, amssymb, amsfonts}
\usepackage{bm}
\usepackage{pgfplots}

% \usepackage{enumitem} <--- Gây lỗi, đã bỏ

\usepackage[
  backend=biber,
  style=authoryear,
  sorting=none,
  bibencoding=utf8
]{biblatex}
\addbibresource{references.bib}

% Fix: "in" bị dính (inNAACL-HLT -> in NAACL-HLT)
\renewbibmacro*{in:}{\printtext{\bibstring{in}\space}}

% 2) Fix: bỏ ngoặc kép quanh title để hết dấu ?
\DeclareFieldFormat[article,inproceedings,book,online,misc]{title}{#1}



\pgfplotsset{compat=1.10}
\PassOptionsToPackage{unicode}{hyperref}
\hypersetup{unicode=true}

\usepackage{xcolor}
\definecolor{myblue}{RGB}{31,119,180}
\definecolor{mygreen}{RGB}{44,160,44}
\definecolor{myorange}{RGB}{255,127,14}
\definecolor{myred}{RGB}{214,39,40}
\definecolor{mypurple}{RGB}{148,103,189}
\definecolor{myblack}{RGB}{0,0,0}

\pgfplotsset{
  cycle list={{myblue},{mygreen},{myorange},{myred},{mypurple},{myblack}},
  every axis plot/.append style={line width=1.2pt, mark=none},
}

\usetheme{Madrid}

% --- CẤU HÌNH ĐÁNH SỐ 1. 2. 3. CHUẨN BEAMER ---
% Thay vì dùng setlist, ta dùng lệnh này của Beamer:
\setbeamertemplate{enumerate items}[default] 

% Hiển thị tên phần ở footer mỗi slide
\setbeamertemplate{footline}{
  \leavevmode%
  \hbox{%
    \begin{beamercolorbox}[wd=.4\paperwidth,ht=2.25ex,dp=1ex,center]{author in head/foot}%
      \usebeamerfont{author in head/foot}\insertsection
    \end{beamercolorbox}%
    \begin{beamercolorbox}[wd=.3\paperwidth,ht=2.25ex,dp=1ex,center]{title in head/foot}%
      \usebeamerfont{title in head/foot}\insertshortauthor
    \end{beamercolorbox}%
    \begin{beamercolorbox}[wd=.3\paperwidth,ht=2.25ex,dp=1ex,right]{date in head/foot}%
      \usebeamerfont{date in head/foot}\insertframenumber{} / \inserttotalframenumber\hspace*{2ex}
    \end{beamercolorbox}}%
  \vskip0pt%
}

% --- THÔNG TIN BÀI THUYẾT TRÌNH ---
\title[Design Patterns]{\large Design Patterns\\[0.2em]
\normalsize trong Phát triển Phần mềm}

\subtitle{\scriptsize Phân tích và vận dụng các\\[0.2em]
\scriptsize Behavioral, Structural và Creational Patterns}

\author[Nhóm 10]{\textbf{Nhóm 10}\\[0.35em]\scriptsize
23127006 -- Trần Nguyễn Khải Luân\\
23127179 -- Nguyễn Bảo Duy\\
23127404 -- Lê Tuấn Lộc}

\institute[HCMUS]{\scriptsize Trường Đại học Khoa học Tự nhiên, ĐHQG-HCM\\[0.25em]
\scriptsize \textit{Học phần:} Thiết kế phần mềm}


\begin{document}

% Title Page
\begin{frame}
  \titlepage 
\end{frame}

% Table of Contents
\begin{frame}{Nội dung trình bày}
  \tableofcontents
\end{frame}

% =========================================================
% SECTION: Refactoring Auction State Logic
% Issues #2, #6, #14 — auction-state.js
% =========================================================
\section{Phân tích \& Refactoring: Auction State Logic}

% ----------------------------------------------------------
% ISSUE #2 — Slide 1: Problem
% ----------------------------------------------------------
\begin{frame}{Issue \#2 — Vi phạm DRY: Duplicate \texttt{productStatus} Logic}
  \textbf{Vấn đề:} Khối logic xác định trạng thái auction được copy nguyên vẹn
  vào 2 route handler khác nhau.

  \bigskip
  \textbf{Nguyên tắc vi phạm:} \alert{DRY} (Don't Repeat Yourself)

  \bigskip
  \begin{block}{Business rule bị encode 2 lần}
    \begin{itemize}
      \item \texttt{GET /detail} → 5-nhánh if/else (dòng $\sim$162)
      \item \texttt{GET /complete-order} → 4-nhánh if/else (dòng $\sim$984)
      \item Hai bản đã \textbf{drift}: bản thứ hai thiếu nhánh \texttt{ACTIVE} cuối
    \end{itemize}
  \end{block}

  \bigskip
  \textbf{Rủi ro kỹ thuật:}
  \begin{itemize}
    \item Thêm state \texttt{DISPUTED} → phải sửa \textbf{cả hai chỗ}; nếu quên một, hai route trả kết quả mâu thuẫn nhau với cùng một product
    \item Không có unit nào để test — logic embedded sâu trong HTTP handler
    \item Business rule bị coupled với HTTP routing layer
  \end{itemize}
\end{frame}

% ----------------------------------------------------------
% ISSUE #2 — Slide 2: Solution
% ----------------------------------------------------------
\begin{frame}[fragile]{Issue \#2 — Refactoring: Extract \texttt{resolveAuctionStatus()}}
  \textbf{Chiến lược:} Extract Function → đưa về \texttt{services/auction/auction-state.js}

  \begin{columns}[t]
    \column{0.48\textwidth}
    \textbf{\alert{Trước}} — inline, lặp lại:
    \footnotesize
    \begin{verbatim}
// Cùng khối xuất hiện 2 lần
const now = new Date();
const end = new Date(product.end_at);
let status = 'ACTIVE';
if (product.is_sold === true)
  status = 'SOLD';
else if (...)
  status = 'CANCELLED';
// ... 3 nhánh nữa
    \end{verbatim}

    \column{0.48\textwidth}
    \textbf{\textcolor{mygreen}{Sau}} — một hàm duy nhất:
    \footnotesize
    \begin{verbatim}
// auction-state.js
export function
  resolveAuctionStatus(product) {
  if (product.is_sold === true)
    return 'SOLD';
  if (product.is_sold === false)
    return 'CANCELLED';
  // ...
  return 'ACTIVE';
}

// Trong route handler:
const status =
  resolveAuctionStatus(product);
    \end{verbatim}
  \end{columns}

  \bigskip
  \textbf{Kết quả:} Pure function, zero dependency, test được bằng \texttt{assert(resolveAuctionStatus(mockProduct) === 'PENDING')}
\end{frame}

% ----------------------------------------------------------
% ISSUE #6 — Slide 1: Problem
% ----------------------------------------------------------
\begin{frame}{Issue \#6 — Vi phạm OCP: State Taxonomy Không Tập Trung}
  \textbf{Vấn đề:} Logic phân loại state auction được viết trực tiếp trong route handlers dưới dạng if/else inline — không có module "đóng" bảo vệ core logic.

  \bigskip
  \textbf{Nguyên tắc vi phạm:} \alert{OCP} (Open/Closed Principle)

  \bigskip
  \begin{alertblock}{Phân biệt quan trọng}
    Guard condition (\texttt{if (!product) return 404}) \textbf{không} phải OCP violation.\\
    State classification (\texttt{if (...) status = 'PENDING'}) \textbf{là} OCP violation — đây là \emph{trục mở rộng domain}.
  \end{alertblock}

  \bigskip
  \textbf{Hậu quả khi thêm state \texttt{PAUSED}:}
  \begin{itemize}
    \item Phải \textbf{mở (modify)} code đang chạy trong production ở nhiều file
    \item Không có "điểm mở rộng" định nghĩa sẵn — developer phải grep toàn codebase
    \item Mỗi file có thể implement logic mở rộng theo cách khác nhau → inconsistency
  \end{itemize}
\end{frame}

% ----------------------------------------------------------
% ISSUE #6 — Slide 2: Solution
% ----------------------------------------------------------
\begin{frame}{Issue \#6 — Refactoring: Encapsulate Variation (State Module)}
  \textbf{Chiến lược:} Tập trung toàn bộ state taxonomy vào một module duy nhất.\\
  Module này \textbf{đóng} với modification, \textbf{mở} cho extension.

  \bigskip
  \begin{block}{\texttt{auction-state.js} — Single Point of Extension}
    \begin{itemize}
      \item \texttt{resolveAuctionStatus(product)} → trả về state string
      \item \texttt{canBid(status)} → guard function
      \item \texttt{canAccessOrder(status)} → guard function
      \item \texttt{requiresAccessCheck(status)} → guard function
    \end{itemize}
  \end{block}

  \bigskip
  \textbf{Thêm state \texttt{PAUSED} sau refactor:}
  \begin{enumerate}
    \item Thêm 1 điều kiện vào \texttt{resolveAuctionStatus()} trong \texttt{auction-state.js}
    \item Thêm guard function \texttt{isPaused(status)} nếu cần
    \item \textbf{Không sửa} bất kỳ route handler hay notifier nào
  \end{enumerate}

  \bigskip
  \textbf{Pattern áp dụng:} \emph{Encapsulate Variation} — mọi variation của "state" sống trong một module. Caller chỉ nhận kết quả, không biết logic phân loại.
\end{frame}

% ----------------------------------------------------------
% ISSUE #14 — Slide 1: Problem
% ----------------------------------------------------------
\begin{frame}{Issue \#14 — Vi phạm DRY: \texttt{auctionEndNotifier} Tự Tính State}
  \textbf{Vấn đề:} \texttt{auctionEndNotifier.js} — module chạy độc lập — tự encode business rule về auction state theo cách riêng, tạo ra \textbf{nguồn sự thật thứ ba}.

  \bigskip
  \textbf{Nguyên tắc vi phạm:} \alert{DRY} ở cấp \emph{cross-module knowledge duplication}

  \bigskip
  \begin{block}{Ba nguồn sự thật trong cùng một hệ thống}
    \begin{itemize}
      \item \texttt{GET /detail}: 5-nhánh if/else (kiểm tra \texttt{is\_sold} + \texttt{highest\_bidder\_id})
      \item \texttt{GET /complete-order}: 4-nhánh if/else (đã drift)
      \item \texttt{auctionEndNotifier.js}: \texttt{if (auction.highest\_bidder\_id)} — 2-nhánh, \textbf{không kiểm tra \texttt{is\_sold}}
    \end{itemize}
  \end{block}

  \bigskip
  \textbf{Bug tiềm ẩn:} Product đã SOLD vẫn có \texttt{highest\_bidder\_id} $\neq$ null.\\
  Nếu query trả về SOLD product, notifier sẽ gửi email ``Bạn đã thắng đấu giá!'' cho một giao dịch đã hoàn tất từ trước.
\end{frame}

% ----------------------------------------------------------
% ISSUE #14 — Slide 2: Solution
% ----------------------------------------------------------
\begin{frame}[fragile]{Issue \#14 — Refactoring: Notifier Depend vào \texttt{auction-state.js}}
  \textbf{Chiến lược:} Notifier không tự tính state — import và dùng \texttt{resolveAuctionStatus()}.

  \begin{columns}[t]
    \column{0.46\textwidth}
    \textbf{\alert{Trước}} — implicit, thiếu check:
    \footnotesize
    \begin{verbatim}
// implicit state check, 2 nhánh
// không check is_sold
if (auction.highest_bidder_id) {
  await sendMail(winner, ...);
  await sendMail(seller, ...);
} else {
  await sendMail(seller, ...);
}
    \end{verbatim}

    \column{0.50\textwidth}
    \textbf{\textcolor{mygreen}{Sau}} — explicit, đúng rule:
    \footnotesize
    \begin{verbatim}
import { resolveAuctionStatus }
  from '../services/auction/
         auction-state.js';

const status =
  resolveAuctionStatus(auction);

if (status === 'PENDING') {
  await notifyWinner(auction);
  await notifySeller(auction);
} else if (status === 'EXPIRED') {
  await notifySellerNoBid(auction);
}
// SOLD/CANCELLED: không gửi email
    \end{verbatim}
  \end{columns}

  \bigskip
  \textbf{Key Takeaway:}
  \begin{itemize}
    \item \textbf{Single Source of Truth}: thay đổi rule tại 1 nơi, propagate toàn hệ thống
    \item \textbf{Dependency on Abstraction}: notifier depend vào function contract, không depend vào raw DB fields
  \end{itemize}
\end{frame}

% ----------------------------------------------------------
% Summary Slide
% ----------------------------------------------------------
\begin{frame}{Tổng Kết: Ba Issues, Một Giải Pháp}
  \begin{block}{Root Cause chung}
    Không có \emph{single source of truth} cho auction state logic.\\
    Issues \#2, \#6, \#14 là cùng một vấn đề nhìn từ ba góc độ khác nhau.
  \end{block}

  \bigskip
  \begin{columns}[t]
    \column{0.48\textwidth}
    \textbf{Trước refactor:}
    \begin{itemize}
      \item 3 nơi encode cùng 1 rule
      \item Drift đã xảy ra giữa các bản sao
      \item Không thể unit test
      \item Thêm state → sửa $\geq 3$ file
    \end{itemize}

    \column{0.48\textwidth}
    \textbf{Sau refactor:}
    \begin{itemize}
      \item 1 module: \texttt{auction-state.js}
      \item Pure function, zero infrastructure dependency
      \item Test: \texttt{resolveAuctionStatus(mock)} trực tiếp
      \item Thêm state → sửa \textbf{1 file}
    \end{itemize}
  \end{columns}

  \bigskip
  \begin{center}
    \textit{``Don't repeat yourself'' không chỉ về code — về \textbf{knowledge}: mỗi business rule chỉ nên tồn tại tại một nơi duy nhất trong hệ thống.}
  \end{center}
\end{frame}

\end{document}